\chapter{Conformance at Scale: A Standalone ACATS Sweep}
\label{ch:acats-sweep}

The runtime is only as trustworthy as the conformance evidence behind it. This
chapter describes how the project runs the official ACATS 4.2 suite \emph{one test
per image} across a pool of ESP32-S3 boards in parallel, and why that arrangement
recovers hundreds of passes that a bundled sweep loses.

\medskip
\noindent\fbox{\parbox{\dimexpr\linewidth-2\fboxsep-2\fboxrule\relax}{\small
\textbf{Note.} The ACATS test suite and its sweep harness
(\texttt{acats\_build.py}, \texttt{acats\_run.py}, the \texttt{ACATS/} tree) are
not shipped in this distribution; they belong to the development repository.
This chapter is kept as a record of the conformance methodology and results; the
commands below are shown to document \emph{how} the sweep was run, not as steps to
reproduce from this repository.}}
\medskip

\section{Two ways to run the suite}\index{ACATS}\index{conformance testing}\index{test harness}

The earlier sweep (\texttt{parallel\_sweep.py}) is \emph{bundled}: dozens of tests
share one flashed image and advance through it with a panic-resume cursor, so a
crash or hang resets the board and the next boot continues past the offending test.
That is fast and cheap on flash cycles, but it couples the tests. A neighbour's
heap churn or secondary-stack residue can exhaust memory mid-batch and reset a later
test \emph{before} it prints its grade: the test is recorded as ungraded even
though it passes in isolation. On the \texttt{full} profile the bundled run grades
914 PASS, with a large ungraded tail dominated by these batch artifacts.

The \emph{standalone} sweep removes the coupling completely: every test is compiled
as its own single-test program, flashed alone, run alone, and graded alone. There is
no shared cursor, no shared heap, no neighbour. The cost (roughly 1{,}500 separate
cross-compiles and flashes) is absorbed by running the work across many boards and
many build trees at once.

\section{Getting the grade off the chip}\index{USB-Serial-JTAG}\index{esp\_rom\_printf}

A subtlety decides whether the standalone approach is even possible. The bare runtime
sends its console (including the ACATS \texttt{Report} \texttt{PASSED}/\texttt{FAILED}
lines) out the ESP32-S3's \emph{built-in USB-Serial-JTAG}, not the UART pins. The
board pool here is wired with ordinary external USB-serial adapters (FTDI and CH343)
on UART0, which never see that JTAG console.

The fix is small. \texttt{ACATS/target/report.adb} routes each grade line through the
ROM's \texttt{esp\_rom\_printf} in addition to \texttt{Ada.Text\_IO}:

\begin{lstlisting}
procedure Line (Msg : String) is
   procedure Rom_Printf (Fmt, Arg : System.Address)
     with Import, Convention => C, External_Name => "esp_rom_printf";
   Fmt : constant String := "%s" & ASCII.LF & ASCII.NUL;
   Buf : String (1 .. Msg'Length + 1);
begin
   Buf (1 .. Msg'Length) := Msg;
   Buf (Buf'Last) := ASCII.NUL;
   Rom_Printf (Fmt'Address, Buf'Address);
end Line;
\end{lstlisting}

The ROM console reaches UART0, so the grades land on the external bridge regardless
of how the board is wired.

\section{The grader must not trust the test id}\index{grading}

The first version of the harness mis-scored most tests as hung. The cause was the
grading regex: it looked for the list name followed by \texttt{PASSED}, but the list
says \texttt{A22006} while \texttt{Report} prints the \emph{unit} name \texttt{A22006B}
-- the suffix letter breaks an exact-id match. Because each image contains exactly one
test, the grader instead matches the \texttt{Report} terminal line \emph{structurally},
ignoring the name:

\begin{lstlisting}[style=sh]
FAILED          ->  **** <anything> FAILED          (checked first)
NOT-APPLICABLE  ->  ++++ <anything> NOT-APPLICABLE
PASSED          ->  ==== <anything> PASSED
TENTATIVELY     ->  !!!! <anything> TENTATIVELY PASSED
\end{lstlisting}

A second early bug scored passing tests as crashes: the flash step reset the board to
run, then the capture step reset it again, producing two \texttt{ESP-ROM:} banners --
and the crash heuristic keyed on ``two or more boot banners.'' Flashing with
\texttt{-{}-no-reset} leaves the chip halted so the capture's own reset is the single
deterministic boot; a second banner then genuinely means a crash-reboot.

\section{Phase 1: parallel build}\index{parallel build}\index{mini-repo}

\texttt{acats\_build.py} compiles each test as a one-test image. The build is the
bottleneck (a single board sits idle for tens of seconds per cross-compile), so it is
parallelised. The complication is that the build performs several \emph{shared} writes:
it regenerates the runtime (\texttt{gen\_runtime.sh}), rebuilds
\texttt{common/bare/boot/obj}, and stages test sources under \texttt{ACATS/work/run},
so naive parallel builds in one tree corrupt each other.

Each builder therefore gets its own isolated ``mini-repo'': a copy of
\texttt{crates}, \texttt{libs}, \texttt{ACATS}, \texttt{ACATS42}, and the two needed
\texttt{examples} subtrees, laid out at the same depth so the build scripts'
\texttt{REPO=\$EX/../..} arithmetic still resolves. The copy carries the warm caches,
so per-test builds stay incremental: no cold runtime rebuild per worker.

\begin{lstlisting}[style=sh]
# OUT/<name>.bin       built OK (the flashable image)
# OUT/<name>.nobuild   could not build / all tests dropped (empty marker)
# OUT/BUILD_DONE        written once every test has been attempted
./acats_build.py ACATS/work/full_applicable.txt OUT 10
\end{lstlisting}

A test that drops to zero buildable units (it needs a library unit the bare runtime
omits) produces \texttt{tests in image: 0} and becomes a \texttt{.nobuild}. The
\texttt{.bin} is written atomically (a temp file renamed into place) so Phase 2 never
sees a partial image, and the phase is resumable: a test that already has a
\texttt{.bin} or \texttt{.nobuild} is skipped.

\section{Phase 2: parallel run and grade}\index{board pool}

\texttt{acats\_run.py} starts one worker per board. Each worker claims the next
un-run \texttt{.bin}, flashes its board with \texttt{-{}-no-reset}, pulses a reset,
captures the console until it sees a grade (or a timeout), classifies the run, and
writes a single marker file holding the raw log:

\begin{lstlisting}[style=sh]
# OUT/<name>.passed  .failed  .na  .crashed  .hung  .nooutput  .flasherr
./acats_run.py OUT 20        # 20-second capture window
\end{lstlisting}

The phases may overlap: Phase 2 drains \texttt{.bin} files as they appear and exits
once \texttt{BUILD\_DONE} exists and no un-run image remains. The whole sweep is
resumable at the file level: delete a marker and re-run, and only that test is
re-flashed. A handful of tests that briefly produced no console (a capture race during
parallel flashing) and several genuinely slow tasking tests were re-run with a 60-second
window simply by removing their markers.

\section{Results}

Full profile, \texttt{full\_applicable.txt} (1{,}518 tests), one image each:

\begin{longtable}{p{4.5cm} r p{6.5cm}}
\toprule
\textbf{Outcome} & \textbf{Count} & \textbf{Meaning} \\ \midrule
\endhead
PASS                & 1{,}286 & graded \texttt{PASSED} / \texttt{TENTATIVELY PASSED} \\
NOT-APPLICABLE      & 10      & test ruled itself out (correctly) \\
FAIL                & 0       & \\
HUNG                & 4       & documented limitations (see below) \\
CRASH               & 0       & \\
won't-build         & 218     & needs a library unit the bare runtime omits \\
\bottomrule
\end{longtable}

That is \textbf{1{,}296 conformant} (PASS plus N/A), \textbf{zero failures}, and
\textbf{zero crashes}, and \textbf{+372 PASS} over the bundled \texttt{full} run
(914), recovered purely by removing batch interaction. Building each test alone is what
turns a co-bundled ``ungraded'' tail back into passes.

\textbf{Three apparent FAIL were harness artifacts, now resolved.} \texttt{CXD1002} and
\texttt{CXD1007} check the \emph{main task's} \texttt{pragma Priority}, live only when the
test is built \emph{as the program main}; the harness now honours
\texttt{standalone\_tests.txt}'s \texttt{main} mode and they PASS. \texttt{CZ1101} is the
ACATS \emph{Report self-test}, which deliberately prints \texttt{**** FAIL\_TEST FAILED}
sub-test sentinels; the grader now matches the result line for the test's \emph{own} name
and takes its final verdict (\texttt{!!!! CZ1101A TENTATIVELY PASSED}).

\textbf{Five apparent HUNG were under-resourced, now passing.} \texttt{CXD1008} and
\texttt{CXD4007} are silent slow tests (\texttt{ImpDef.Switch\_To\_New\_Task} is 1.0\,s, so
they need minutes, not the default 60\,s window); \texttt{C433A01/02} and \texttt{C953002}
OOM'd at the 64\,KB secondary-stack default. The harness now honours two tuning lists:
\texttt{long\_wdt\_tests.txt} (a 240\,s capture window) and \texttt{small\_sec\_stack\_tests.txt}
(the 4\,KB default for many-task tests), and all five PASS.

\textbf{Two more were program-termination cases, now fixed.}
\texttt{C761001}'s grade is emitted from a \emph{library-level} controlled object's
\texttt{Finalize}, which runs only at program termination, but
\texttt{Configurable\_Run\_Time} makes gnatbind omit the \texttt{finalize\_library} chain,
so it never ran. Flipping that flag is not viable (it needs a full runtime rebuild and is
entangled with the runtime's codegen; the \texttt{.ali} flags are identical native-vs-target,
so it is purely the binder mode). Since the per-unit finalizers
(\texttt{<unit>\_\_finalize\_spec/body}) are still generated, the build synthesizes a
finalizer-runner that calls the user units' finalizers in reverse elaboration order (parsed
from the binder's \texttt{adainit}; weak + guarded), merges it into the image, and the bare
env body calls it after \texttt{adafinal}, reached only by a program-main image after its
main returns, so the wrapped sweep images are unaffected. \texttt{C761001} now PASSES.

\texttt{C94004} is a multi-file test whose real main is \texttt{C94004A}, a library task
``examined to see whether it terminates.'' As a wrapped batch test its non-terminating
library tasks pile up and hang; built as the program main, only \texttt{C94004A} runs and
prints \texttt{PASSED} from its own body. Its main unit (\texttt{c94004a}) differs from the
list key (\texttt{C94004}), so the build resolves the boot entry from
\texttt{batch\_runnable.txt} (\texttt{build\_ada.sh} writes \texttt{.noidf/ada\_main.sym},
which \texttt{bare\_build} honours instead of the guessed \texttt{\_ada\_c94004}).
\texttt{C94004} now PASSES.

\textbf{The remaining four HUNG are documented limitations.} \texttt{CXD1006} (an
\texttt{Interrupt\_Priority} task collides with the kernel's own Xtensa level 5),
\texttt{CB1010} (a large-frame recursive stack overflow the 64-byte watchpoint cannot
catch), \texttt{C480001} (unbounded recursion / GNAT codegen), and \texttt{CXC7004}
-- \emph{not} a deadlock but an environment-task \emph{termination} test (JTAG dual-core
post-mortem shows cpu0 spun past \texttt{ADA\_MAIN}, cpu1 idle) whose grade depends on
shutdown-path behavior the Configurable runtime does not provide.

\textbf{The ten N/A are correct.} \texttt{CXD2001-2004} and \texttt{CXD6001-6003} test
\emph{uniprocessor} dispatching (a single ready queue with deterministic ordering),
which does not apply to a genuine two-core SMP target, so each prints
\texttt{+ Multi-Processor Configuration} and grades NOT-APPLICABLE. \texttt{C45322},
\texttt{C45523}, and \texttt{C4A012} rule themselves out for their own reasons.

\subsection{A wider net: the maximal list across seven boards}\index{Ada.Interrupts}\index{Ada.Calendar}

A later run pushed the net wider on purpose. Rather than the curated
\texttt{full\_applicable.txt}, it swept the \emph{maximal union} (the applicable
list plus everything previously parked in \texttt{jorvik\_skipped.txt} and
\texttt{ungraded\_tests.txt}, 1{,}536 tests): one image each, across \emph{seven}
boards at once (the standalone harness is one worker per serial console, so seven
consoles run seven tests in parallel). The point of including the already-skipped
set is that capabilities added since the lists were drawn (the Calendar children,
and the interrupt work) might turn an old \texttt{nobuild} into a pass; running the
superset catches that automatically.

It did, and it also earned its keep by finding a real bug. The dynamic
\texttt{Ada.Interrupts} tests \texttt{CXC3007}/\texttt{CXC3008} \emph{failed},
which is how the restricted-\texttt{a-interr}-stub problem (every dynamic operation
\texttt{raise Program\_Error}; see the Interrupts chapter) came to light. With the
delegating \texttt{a-interr} body in place, \texttt{CXC3007} now \texttt{PASSES} and
\texttt{CXC3008} advances into its interactive ``generate the interrupt'' phase
(which the unattended rig cannot satisfy). The wider net found a second real bug
too: \texttt{C954025}/\texttt{C954026} (a requeue-with-abort whose timeout is a
\texttt{delay until} an \texttt{Ada.Calendar} time) expired at the wrong moment,
because the kernel's timed-sleep collapsed the \texttt{Absolute\_Calendar} and
\texttt{Absolute\_RT} delay modes and fed a Calendar-base deadline to the raw
counter. Shifting a Calendar deadline back into the raw base (adding the BB-time
\texttt{Epoch}) in \texttt{Timed\_Sleep} fixes both. They now \texttt{PASS}
(they are slow and silent, so they join the wide-watchdog list). After these
fixes the residue on the wider list is what you would expect: the documented
frontiers (\texttt{CXD1006}, \texttt{CB1010}, \texttt{C480001}, \texttt{CXC7004})
and the interactive interrupt and suspension tests that need a bench-generated
interrupt. No crashes, and no surprise failures in the conformant core.
(\texttt{C96004}'s post-2099 Calendar range was later closed too, by coarsening the
Julian-day base $5\times$ to reach 2399; see the Limitations chapter. A companion
fix set \texttt{System.Tick} to the true \texttt{1/240\,MHz} clock period (it had
shipped as \texttt{0.0}, telling clock-resolution tests the clock was infinitely
precise), which is what \texttt{C96001}/\texttt{C96004} read.)

\section{The embedded profile}\index{embedded profile}\index{Jorvik}\index{No\_Task\_Hierarchy}

The same harness runs the \texttt{embedded} (ZCX) profile against the
\texttt{jorvik\_hw\_runnable.txt} list (846 tests). It needed one runner fix first. For a
single-test image the generated driver adds a do-nothing task to force the environment
task's secondary stack; that task was declared \emph{nested inside} the \texttt{Main}
procedure, which is fine under \texttt{full} but violates the embedded / light-tasking
Jorvik restriction \texttt{No\_Task\_Hierarchy} (a task nested in a subprogram is illegal).
Every size-1 image therefore failed to build: the standalone sweep produced
all-no-build. The bundled \texttt{parallel\_sweep} never hit this because it pads
isolated images with filler tests to keep the batch larger than one, so the helper task
was never emitted. The fix is to emit the helper as a \emph{library-level} task instead,
which is legal under Jorvik and accepted by \texttt{full} alike.

\begin{longtable}{p{4.5cm} r p{6.5cm}}
\toprule
\textbf{Outcome} & \textbf{Count} & \textbf{Meaning} \\ \midrule
\endhead
PASS                & 840 & graded \texttt{PASSED} / \texttt{TENTATIVELY PASSED} \\
NOT-APPLICABLE      & 3   & test ruled itself out (correctly) \\
FAIL                & 1   & \texttt{CXC3004}, the interactive test (see below) \\
HUNG                & 0   & \\
CRASH               & 0   & \\
won't-build         & 2   & partition-wide build constraints (see below) \\
\bottomrule
\end{longtable}

\textbf{843 conformant, zero HUNG, zero CRASH, zero genuine FAIL}, above the bundled
embedded baseline (836). The one FAIL, \texttt{CXC3004}, is the \emph{interactive} test:
it waits for a manually-generated external interrupt, which nothing raises on the
unattended rig: a documented limitation, not a runtime bug.

The two won't-build are notable because neither is a missing Ada library unit --
each is a \emph{partition-wide} build constraint. \texttt{CXH1001} (Annex H) uses the
configuration pragma \texttt{Normalize\_Scalars}, which must be applied to every unit in
the partition including the runtime; the prebuilt embedded runtime is not compiled with it,
so the binder rejects the program (\emph{``some but not all files compiled with
Normalize\_Scalars''}). \texttt{CA5006} (a \texttt{pragma Elaborate} separate-compilation
test) binds to an \emph{elaboration circularity} under the embedded runtime's static
elaboration model. Both are correctly counted as won't-build rather than failures.

\section{Closing a won't-build category: the Ada.Calendar family}\index{calendar overlay}

Some of the ``won't-build'' tail is not a fundamental limitation but a simply-missing library
unit, and the sweep is what makes that visible. The bare runtime originally shipped only base
\texttt{Ada.Calendar} (plus \texttt{.Conversions}); the RM~9.6.1 \emph{child} units were
absent, so every Calendar formatting/arithmetic test dropped at compile. Those children are now
provided on the \texttt{embedded} and \texttt{full} profiles, from a shared
\texttt{calendar\_overlay/} applied by \texttt{gen\_runtime.sh}:

\begin{itemize}
  \item \texttt{Ada.Calendar.Arithmetic}: day arithmetic on the parent's
        Modified-Julian-Day representation (exact; the bare runtime has no leap-second table,
        so \texttt{Leap\_Seconds} is always zero).
  \item \texttt{Ada.Calendar.Formatting}: \texttt{Image}/\texttt{Value}/\texttt{Split}/
        \texttt{Time\_Of}/\texttt{Day\_Of\_Week}, written against the parent's public API with
        hand-rolled decimal formatting (the runtime excludes the \texttt{'Image}/\texttt{'Value}
        helper units, so the displayed fraction is truncated and \texttt{Value} raises on
        out-of-range fields, both as the language requires).
  \item \texttt{Ada.Calendar.Time\_Zones}: bare metal has no time-zone database, so the
        Calendar keeps UTC and \texttt{Local\_Time\_Offset} returns zero.
  \item \texttt{Ada.Calendar.Delays}: \texttt{delay until} on a \texttt{Calendar.Time},
        delegated to the SMP-hardened \texttt{Ada.Real\_Time.Delays}.
  \item the obsolescent J.1 top-level \texttt{Calendar} rename, so Ada-83 \texttt{with Calendar;}
        tests build.
\end{itemize}

Hardware-verified: \texttt{C961001} (Formatting \texttt{Image}/\texttt{Value} conformance) and
\texttt{C960002} (\texttt{Delay\_Until}) PASS on both profiles, and the \texttt{with Calendar;}
tests \texttt{C96005}/\texttt{C96006} PASS. \texttt{C96004}, which requires
\texttt{Year\_Number'Last = 2399} (Ada~2005), now PASSES as well: rather than the
full time-base rewrite first assumed, the Julian-day base is coarsened $5\times$
(sub-day scale \texttt{Duration'Small\,/\,17\_280.0}), trading 1\,ns Calendar
resolution for 5\,ns and buying the extra three centuries within signed 64 bits --
see the Limitations chapter. \texttt{light-tasking} is deliberately
\emph{excluded}: it is a \texttt{No\_Exception\_Propagation} profile, so every
\texttt{Time\_Error}/\texttt{Constraint\_Error} becomes a board reset. \texttt{Ada.Calendar}
is unusable and grades nothing there, which is exactly why bb-runtimes omits it from the Jorvik
profile. By contrast \texttt{Ada.Numerics} needed no work at all: the parent, elementary
functions, the random generators, complex types, real arrays and big integers already ship and
compile into both profiles (HW-verified via \texttt{CXG2003}, \texttt{CXG2009}, \texttt{CXG1004}).

\section{Why this matters}

The standalone sweep is slower per test, but it produces an unambiguous, file-per-test
record: every outcome is one marker file containing the exact console transcript that
justified it. Cross-checking the two methods is itself a result: the bundled run's
ungraded tail is explained, test by test, by the standalone run's passes, and the only
remaining non-passes are a short, named list of artifacts and documented frontiers.
